% New standalone increment; prior increments and repository are read-only.
\documentclass[11pt]{article}
\usepackage[T1]{fontenc}\usepackage[utf8]{inputenc}\usepackage{lmodern}
\usepackage[margin=.9in]{geometry}
\usepackage{amsmath,amssymb,amsthm,mathtools,microtype,array,longtable,booktabs,xurl}
\usepackage[hidelinks]{hyperref}
\setlength{\emergencystretch}{3em}\allowdisplaybreaks[2]
\numberwithin{equation}{section}
\newtheorem{theorem}{Theorem}[section]
\newtheorem{lemma}[theorem]{Lemma}\newtheorem{proposition}[theorem]{Proposition}
\theoremstyle{remark}\newtheorem{remark}[theorem]{Remark}
\newcommand{\dd}{\,d}\newcommand{\E}{\mathbb E}
\newcommand{\KK}{\mathcal K}\newcommand{\BB}{\mathcal B}
\newcommand{\Gp}{{G_+}}\newcommand{\Clk}{\mathcal C_q}
\newcommand{\eps}{\varepsilon}\newcommand{\one}{\mathbf1}
\newcommand{\Kfun}{\mathsf K}
\newcommand{\src}[2]{\textup{#1} \texttt{#2}}
\title{Matching package inside the initialized weak passage\\
Population accounting, logistic phase and weighted incoming smallness}
\author{Standalone analytical increment for independent review}
\date{Version 1 --- 30 October 2026}
\begin{document}\maketitle
\begin{abstract}
The full clock cohort carries all but a vanishing relative amount of
the non-core mass and the weak response on a late passage interval.
The omitted primary band has relative amplitude of order $q^{1/4}$,
not merely a bounded fraction. The exact amplitude clock then matches
an auxiliary logistic mass with $o(1)$ phase. BULK's separate-clock
transverse theorem applies under RB's proved states and regenerates
$J_S\le C(q^3+q^2H)$, which supplies the previously missing
$\sqrt{N_0}W_2(t_s)\to0$. We also record exact mass identities and
quantitative late core-mass defect budgets. This is a matching theorem
on the original flow, not a weak-field reduction or a selection theorem.
\end{abstract}

\section{Interval, fixed populations and inherited facts}
\label{an03mpv1:setup}
All new labels have prefix \texttt{an03mpv1:}. Work in the initialized
deep regime $\kappa=10$, $\rho\in[13/20,7/10]$, and set
\[
 \lambda=\rho^2,\quad Q=\lambda+q^2,\quad \ell=\log(1/q),
 \quad k=10(1-\lambda),\quad
 \gamma=1-5\lambda+\frac{2\lambda}{1-\lambda}.
\]
Keep WPT/RB's fixed constant order and then WBR's choices.
All passage inputs are conclusions of RB's completed theorem;
we do not reopen its bootstrap. The fixed original label partition is
$\KK,\Gp,\Sigma,\BB$, with $\mathcal D=\varnothing$, where
\begin{align}
 \Gp&=\{b\in[\pi/2,\pi]:c_0=-\cos b\ge Kq^\gamma\},
 &\Sigma&=\{b\in[\pi/2,\pi]:c_0<Kq^\gamma\},\notag\\
 S&=\Clk=\{b\in[\pi/2,\pi]:c_0\ge q^{1/4}\},
 &P_q&=\Gp\setminus S.
 \label{an03mpv1:partition}
\end{align}
Here $K=K(C_0)$ is Q2EC's fixed threshold. The proved bound
$\gamma>7/20$ ensures $Kq^\gamma<q^{1/4}$ for small $q$.
Thus the union called $\Gp$ in LCRC is precisely the primary Q2
population, counted once. Labels are fixed, including the endpoints;
the initialization measure is $\dd\lambda_0=\dd b/(2\pi)$.

Use the exact balanced-flow energies
\begin{gather*}
 z=(U_2+W_2)/2,\quad a_-=(U_2-W_2)/2,\quad
 H_T=\int_Tz^2,\quad A_T=\int_Ta_-^2,\\
 J_T=\tfrac12\int_T(U_1^2+W_1^2),\quad E_T=H_T+A_T.
\end{gather*}
All label integrals use $\lambda_0$. Balance gives $\mu_T=E_T+J_T$.
Write $H=H_S$, $M=\mu_\KK$ for the physical core mass, and
$c_T=Q^{-1}\E_{P_2}[X_2(f_T)_2]$ for the response of a cohort.
Do not identify $M$ with $M_{\rm res}$ or with a normalized mass.

We choose the start exponent
\begin{equation}
 a=1/8,\qquad t_s=\inf\{t\ge t_0:H(t)=q^a\},\qquad
 t_e=\inf\{t\ge t_0:H(t)=h\},
 \label{an03mpv1:levels}
\end{equation}
where $h$ is any fixed level satisfying WBR's smallness conditions;
it may be decreased. In particular $h<c_Hc_b/2$ and $h<1/4$.
RB and WBR give $t_0<t_s<t_e<t_b<t_0+20\ell$ for small $q$.
BULK's clock proof gives $z'/z\ge b_*/2>0$ on $S$, where
$b_*=Q(1-c_b)/2$; hence $H$ is strictly increasing. In particular
$q^a\le H\le h$ on the selected interval.

The following inherited bounds will be used with their stated scopes:
\begin{align}
 &E_\Gp\le C_EH,\quad E_T\le(10/9)H_T\ (T\subseteq\Gp),
 \quad J_G\le q^2K_Y^2,\quad G=\Gp\cup\Sigma,\notag\\
 &E_\Sigma\le C_\Sigma(q^5+q^{\Gamma_\nu}H),\quad
 \Gamma_\nu\ge289169/924000>3/10,\quad
 \mu_\BB\le q^{11/10},\quad |c_\KK|\le Cq^2,
 \label{an03mpv1:states}\\
 &H_0:=H(t_0)\asymp A_q^{-\lambda}q^k,\quad
 0<c_A\le A_q\le C_A,\quad
 \int_s^t|\eps_H|\le D_q:=C(q^{1/8}+q^2\ell),\notag\\
 & (\log H)'=Q(1-c)+\eps_H,\quad 0<c\le c_b\le1/4,
 \quad b_0=Q(1-c_b)>0.
 \label{an03mpv1:clock}
\end{align}
The absolute defect integral is BULK's conclusion, certified by RB,
not a consequence inferred by differentiating a bounded phase error.
It implies
\begin{equation}
 H(v)\le e^{D_q}H(t)e^{-b_0(t-v)},\qquad
 \int_{t_0}^t H\le\frac{e^{D_q}}{b_0}H(t),\quad
 \int_{t_0}^t\sqrt H\le\frac{2e^{D_q/2}}{b_0}\sqrt{H(t)}.
 \label{an03mpv1:backward}
\end{equation}
The weak coordinates are $u=(\pi/2-\theta)/q$ and
$v=(\psi-\pi/2)/q$. WBR gives $u>0$ on $S$ from $t_0$ to $t_e$,
$|\theta-\psi|\le C_\Delta q$, and, on $[t_s,t_e]$,
\begin{equation}
 |u|+|v|\le L(q):=C_L(1+q^{-p_*}),\qquad p_*=603/640<1.
 \label{an03mpv1:chart}
\end{equation}
The reached chart, not a bounded-chart assumption, justifies all
small-angle estimates below.

\section{Step 1: population accounting and signed response}
\begin{lemma}[The omitted primary band]
\label{an03mpv1:band}
Uniformly on the entire RB passage,
\begin{equation}
 H_{P_q}(t)/H(t)\asymp q^{1/4}.
 \label{an03mpv1:bandratio}
\end{equation}
Both constants are fixed and may depend on the primary threshold and
the fixed RCET distortion. In particular this is a vanishing ratio.
\end{lemma}
\begin{proof}
RB \texttt{primaryseed}, together with Q2EC's two-sided original
coordinate transport in its matched primary sector, gives
\begin{equation}
 U_2(t_0,b)^2\asymp A_q^{-\lambda}q^k(s_0+q)^2,
 \qquad s_0=\sin b,\quad b\in\Gp.
 \label{an03mpv1:seed}
\end{equation}
At entry $|a_-|/z\le1/4$, so $z^2\asymp U_2^2$.
The Q2EC matching proof includes late crossings, particularly $b=\pi$,
and retains both $s_0+q$ and $c_0+q$. Its proof also checks the bounded
shift $t_0-10\ell=-\log A_q+o(1)$. Thus this is not a transfer
through an absolute angle error or a logarithm of the initial zero seed.

On $P_q$, $Kq^\gamma\le c_0<q^{1/4}$, so $s_0\asymp1$.
The change of variable $c_0=-\cos b$ has density
$(2\pi\sqrt{1-c_0^2})^{-1}$, bounded above and below there.
Because $\gamma>7/20>1/4$,
\[
 \lambda_0(P_q)\asymp q^{1/4}-Kq^\gamma\asymp q^{1/4}.
\]
Integrating \eqref{an03mpv1:seed} gives
$H_{P_q}(t_0)\asymp A_q^{-\lambda}q^{k+1/4}$, whereas
$H(t_0)\asymp A_q^{-\lambda}q^k$.
RCET \texttt{transport} applies to the two fixed positive-amplitude
subsets $P_q,S\subseteq\Gp$, with all its response, duration, residual,
learning and Cartesian entry hypotheses discharged by RB/Q2EC.
It multiplies their relative ratio by a factor in $[e^{-2D_+},e^{2D_+}]$,
where $D_+=O(C_0+q^2\ell)$ is bounded. This proves
\eqref{an03mpv1:bandratio}; it does not claim $D_+=o(1)$.
\end{proof}

\begin{proposition}[Accounting on the late interval]
\label{an03mpv1:accounting}
Uniformly for $t_s\le t\le t_e$,
\begin{align}
 H_{P_q}+\mu_\Sigma+\mu_\BB&\le Cq^{1/4}H,
 \label{an03mpv1:omitted}\\
 |c-H|&\le Cq^{1/4}H+Cq^2.
 \label{an03mpv1:response}
\end{align}
Consequently the omitted weak mass is $o(H)$ even if all outside
mass, rather than only its weak-side part, is charged. No population
mask is differentiated.
\end{proposition}
\begin{proof}
The secondary mass satisfies
\[
 \frac{\mu_\Sigma}{H}\le C\left(q^{\Gamma_\nu}
       +q^{5-a}+q^{2-a}\right),\qquad
 \frac{\mu_\BB}{H}\le q^{11/10-a}.
\]
For $a=1/8$ the latter exponent is $39/40$, and all the displayed
exponents exceed $1/4$. Together with Lemma \ref{an03mpv1:band}
this proves \eqref{an03mpv1:omitted}.

To estimate the selected response, WBR alignment and balance give
\begin{equation}
 A_S\le(C_\Delta^2/4)q^2\mu_S,\qquad
 \mu_S=H+A_S+J_S\le C(H+q^2)\le CH.
 \label{an03mpv1:earlymass}
\end{equation}
The last step uses $H\ge q^a$. On the reached chart
$\sin\theta\ge1/2$ for small $q$. In LCRC's exact ledger,
\[
 c_S=H-A_S+\tau_S,\qquad
 \tau_S=Q^{-1}\int_S W_2\E_{P_2}[X_2(U\cdot X)_-],
\]
Gaussian projection gives the two-sided absolute estimate
\begin{equation}
 |\tau_S|\le C\mu_S e^{-c_*/q^2}.
 \label{an03mpv1:tail}
\end{equation}
Indeed $|W_2||U|\le m$, the gate-failure probability is at most
$\exp[-\rho^2/(8q^2)]$, and Cauchy--Schwarz bounds the remaining
fourth moment uniformly. This uses the reached input chart; LCRC's
one-sided \texttt{taulower} alone would not give an upper error.
Thus $|c_S-H|\le Cq^2H+CH e^{-c_*/q^2}$.

For each other cohort LCRC \texttt{responsecharge} and Young's
inequality give
\[
 |c_T|\le E_T+\frac{2q}{\sqrt Q}\sqrt{E_TJ_T}
          \le2E_T+\frac{q^2}{Q}J_T.
\]
Hence $|c_{P_q}|\le Cq^{1/4}H+Cq^4$ and
$|c_\Sigma|\le C(q^5+q^{\Gamma_\nu}H+q^4)$.
The core response is $O(q^2)$, and
$|c_\BB|\le C\mu_\BB\le Cq^{39/40}H$.
Sum the exact decomposition
$c=c_S+c_{P_q}+c_\Sigma+c_\KK+c_\BB$ to obtain
\eqref{an03mpv1:response}. All four data/gate contributions remain
in the exact ledger; no sign of an omitted cohort is assumed.
\end{proof}

\section{Step 2: logistic matching with a vanishing phase error}
Define the auxiliary scalar $N$ by
\begin{equation}
 N'=\lambda N(1-N),\qquad N(t_s)=N_0=H(t_s)=q^a.
 \label{an03mpv1:logistic}
\end{equation}
It is a comparison mass, not a declaration of the original radial law.

\begin{proposition}[Uniform phase and interior hitting times]
\label{an03mpv1:phase}
On $[t_s,t_e]$, with $x=\log(H/N)$,
\begin{equation}
 \sup|x|\le d_{\rm ph}(q):=
 C\left(q^{1/8}+q^2\ell+\frac{q^{1/4}h}{b_0}\right)=o(1).
 \label{an03mpv1:phasebound}
\end{equation}
In particular $N(t_e)=h\exp(O(d_{\rm ph}))<2h<1/2$ for small $q$.
For every fixed $n\in(0,h)$, both first hits $t_H(n),t_N(n)$ lie
in $[t_s,t_e]$ for small $q$, and
\begin{equation}
 |t_H(n)-t_N(n)|\le C\frac{d_{\rm ph}}{\lambda(1-h)}=o(1).
 \label{an03mpv1:hittime}
\end{equation}
For $n=h$, the same $o(1)$ relation holds to the auxiliary logistic
hit, which may lie just after $t_e$; no original-flow estimate after
$t_e$ is used in that endpoint comparison.
\end{proposition}
\begin{proof}
The exact identity is
\[
 x(t)=q^2\int_{t_s}^t(1-c)+\lambda\int_{t_s}^t(N-c)
                   +\int_{t_s}^t\eps_H.
\]
Its derivative can be reorganized as
\begin{equation}
 x'=\lambda N(1-e^x)+f(t),\qquad
 f=q^2(1-c)+\lambda(H-c)+\eps_H.
 \label{an03mpv1:contractphase}
\end{equation}
For every real $x$, the first term opposes its sign, so
$D^+|x|\le|f|$. Equations \eqref{an03mpv1:backward} and
\eqref{an03mpv1:response} give
\[
 \int_{t_s}^{t_e}|H-c|\le Cq^{1/4}\frac{e^{D_q}h}{b_0}
                                  +Cq^2\ell.
\]
Using the absolute clock-defect integral proves
\eqref{an03mpv1:phasebound}, without a factor $e^{Ch\ell}$.
As an additional useful conclusion,
\begin{equation}
 \int_{t_s}^{t_e}|N-c|
 \le C\left(d_{\rm ph}\frac{h}{b_0}
              +q^{1/4}\frac{h}{b_0}+q^2\ell\right)=o(1).
 \label{an03mpv1:phaseintegral}
\end{equation}
Here $|N-H|\le C d_{\rm ph}H$ follows from the uniform bound.

At $t_H(n)$, $N=n e^{-x(t_H(n))}$. Exact logistic integration
therefore gives
\[
 t_H(n)-t_N(n)=\lambda^{-1}
 \left[-x(t_H(n))+
 \log\frac{1-n}{1-n e^{-x(t_H(n))}}\right].
\]
The mean value theorem proves \eqref{an03mpv1:hittime}. For fixed
$n<h$, $N(t_e)>n$ for small $q$, so the logistic hit is interior.
At $n=h$ the formula uses only the explicit auxiliary logistic
solution, not an extension of any original estimate.
\end{proof}

\section{Step 3: mass identities, without identifying the three masses}
\begin{proposition}[Selected and whole weak mass]
\label{an03mpv1:masses}
The exact identities and current pointwise consequences are
\begin{align}
 N_S:=\mu_S&=H+A_S+J_S,\qquad A_S\le Cq^2N_S,
 \label{an03mpv1:identity}\\
 \frac{N_S}{H}&=1+O(q^2+q^{2-a})=1+O(q^{15/8}).
 \label{an03mpv1:coarsemass}
\end{align}
The transverse estimate proved in the next section sharpens this to
\begin{equation}
 \frac{N_S}{H}=1+O\left(q^2+\frac{q^3}{H}\right)=1+O(q^2).
 \label{an03mpv1:sharpmass}
\end{equation}
Define the full fixed non-core radial mass by
$N_{\rm nc}=\mu_\Gp+\mu_\Sigma+\mu_\BB$. Then
\begin{equation}
 \frac{N_{\rm nc}}H=1+O(q^{1/4}),\qquad
 \frac{N_S}N=1+O(d_{\rm ph}+q^2),\qquad
 \frac{N_{\rm nc}}N=1+O(d_{\rm ph}+q^{1/4}).
 \label{an03mpv1:totalmass}
\end{equation}
Any weak-family mass containing $S$ and drawn from these non-core
cohorts obeys the same $1+O(q^{1/4})$ comparison with $H$.
If only part of $\BB$ is called weak, charging all of $\BB$ proves
the assertion without needing a moving-mask evolution law.
\end{proposition}
\begin{proof}
Balance proves the identity. Since $|\sin\theta-\sin\psi|
\le|\theta-\psi|\le C_\Delta q$, WBR alignment proves the bound
for $A_S$. Absorb $Cq^2N_S$ and use $J_S\le q^2K_Y^2$
to get \eqref{an03mpv1:coarsemass}. Inserting the next section's
$J_S\le C(q^3+q^2H)$ instead gives \eqref{an03mpv1:sharpmass}.
For the omitted primary group,
$\mu_{P_q}\le(10/9)H_{P_q}+J_{P_q}$, with
$J_{P_q}\le q^2K_Y^2$. Add Proposition
\ref{an03mpv1:accounting} and use $q^2/H\le q^{15/8}$.
The ratios to $N$ then follow from the phase estimate.
No circular use of \eqref{an03mpv1:sharpmass} was made in that
phase proof or in the transverse input below.
\end{proof}

\section{Step 4: regenerated transverse energy and second-moment smallness}
\begin{lemma}[The proved transverse theorem applies]
\label{an03mpv1:transverse}
There is a fixed $C_J'$ such that, on the entire certified passage,
\begin{equation}
 J_S(t)\le C_J'[q^3+q^2H(t)].
 \label{an03mpv1:J}
\end{equation}
\end{lemma}
\begin{proof}
Apply BULK \texttt{forcing} with its cohort $G=S$ and exponent
$\sigma=0$. Its complete inputs are as follows:
the fixed positive clock $H$, $0\le c\le c_b<1$, the subinterval
lower bound $\int\eps_H\ge-D_q$, logarithmic duration,
$E_S\le C_E(q^5+H)$, $J_S(t_0)\le C_Jq^3$,
$\ell_S\le L_\times q$, and $\int r_S\le B$.
The clock and response range are \eqref{an03mpv1:clock};
signgood gives $E_S\le10H/9$, which is independent of
\eqref{an03mpv1:J}. Q2EP \texttt{entryenergy} applies to every
fixed Q2 subcohort and gives the required initial $J_S$.
RB's residual norms imply $\ell_S\le L_\times q$ by BULK
\texttt{residualinputs}; RB also supplies $\int r_S\le B$
and duration at most $20\ell$. All constants are fixed, and
$D_q$ can be bounded by one.
Thus BULK \texttt{Jprop} gives \eqref{an03mpv1:J} directly.
Equivalently its already proved norm inequality and
\eqref{an03mpv1:backward} bound $\sqrt{J_S}$ by a fixed multiple
of $q^{3/2}+q\sqrt H$. No leading-system capture or bounded
weak chart is a hypothesis of this Cartesian theorem.
\end{proof}

For initial comparison, set
\[
 \eta_s(\dd b)=\frac{m(t_s,b)}{N_S(t_s)}\dd\lambda_0(b)
 \quad(b\in S),\qquad N_0=q^a.
\]
This is a positive probability on fixed labels. Let
$(U_*,V_*)(N)$ be the selected leading branch and define, just at
entry,
\begin{equation}
 W_2(t_s)=\big\||u-U_*(N_0)|+|v-V_*(N_0)|\big\|_{L^2(\eta_s)}.
 \label{an03mpv1:Wp}
\end{equation}
AN06 \texttt{regiontheorem}, as inventoried in SI, gives
$0<U_*<1$, $0<V_*<11/10$ on $0<N\le1/2$ in the ratio box.
This is used solely for bounded comparator coordinates, not as
an original-field or local-strong-guard assertion.

\begin{proposition}[Weighted incoming smallness]
\label{an03mpv1:smallness}
With $a=1/8$ and $p=2$,
\begin{equation}
 N_0=q^{1/8}\longrightarrow0,\qquad
 A_0:=\sqrt{N_0}W_2(t_s)
 \le C(q^{1/2}+q^{1/16})=O(q^{1/16})\longrightarrow0.
 \label{an03mpv1:A0}
\end{equation}
The same bound holds if the initial scalar normalization is chosen
as the actual $N_S(t_s)$ instead of $H(t_s)$, with a corresponding
bounded comparator at that mass. In that convention the initial
logarithmic phase is $O(q^2)$ rather than zero.
\end{proposition}
\begin{proof}
On the reached chart, $|qu|,|qv|\le1/2$ for small $q$, so
\[
 \int_S m(u^2+v^2)\le C J_S/q^2.
\]
Minkowski and the bounded comparator give
\[
 N_0W_2(t_s)^2\le C\frac{N_0}{N_S(t_s)}
             \left[\frac{J_S(t_s)}{q^2}+N_S(t_s)\right].
\]
Use \eqref{an03mpv1:sharpmass} and \eqref{an03mpv1:J} at
$H(t_s)=q^a$ to bound the right side by $C(q+q^a)$.
This proves \eqref{an03mpv1:A0}. The changed normalization assertion
uses $N_S(t_s)/N_0=1+O(q^2)$ and the same uniform comparator bounds.
No assumption of constant future normalized original weights is made.
\end{proof}

\paragraph{Audit of the proposed supremum route.}
At a first clock level $H=q^{\widetilde a}$, SI's BULK bound gives
the exponent
\[
 p(\widetilde a)=15/8-(k-\widetilde a)(1-c_b)/4.
\]
For this estimate alone to prove both a reached small-angle chart
and $\sqrt{H}\sup(|u|+|v|)\to0$, the strict inequalities are
$p(\widetilde a)<1$ and $\widetilde a/2-p(\widetilde a)>0$.
Equivalently,
\begin{equation}
 \frac{15/2-k(1-c_b)}{1+c_b}<\widetilde a
       <k-\frac{7}{2(1-c_b)}.
 \label{an03mpv1:supinterval}
\end{equation}
On this box such an interval can exist only if
\begin{equation}
 k>\frac{11}{2}+\frac{7c_b}{2(1-c_b)}.
 \label{an03mpv1:supthreshold}
\end{equation}
Thus the proposed $k>5.5$ count is a necessary limiting condition,
and at fixed positive $c_b$ the actual condition is stricter.
It fails on part of the box, since $k=51/10$ at $\rho=7/10$.
This failure of a coarse sufficient bound is not an obstruction to
\eqref{an03mpv1:A0}; the proved transverse theorem supplies the
needed moment estimate throughout the box.

\section{Step 5: late physical-core mass-defect budgets}
WBR \texttt{masseq, massrate} gives on the completed passage
\begin{equation}
 (\log M)'=1-M+\eta,\quad |\eta|\le C_\eta q^{11/10},\qquad
 |1-M(t)|\le C_me^{-(t-t_0)/3}+6C_\eta q^{11/10}.
 \label{an03mpv1:massrate}
\end{equation}
RB's exact clock and bounded $\int c$ imply
\begin{equation}
 t_s-t_0=\frac{(k-a)\ell+\lambda\log A_q+O(1)}Q,
 \quad \alpha_M:=\frac{k-a}{3Q}\ge\frac{199}{60}
 \label{an03mpv1:starttime}
\end{equation}
for small $q$, using $Q\le1/2$ and $k\ge51/10$.
Thus, with $w=t-t_s\ge0$,
\[
 |1-M(t)|\le C[q^{\alpha_M}e^{-w/3}+q^{11/10}].
\]
BULK's size can also be retained with its time decay:
\[
 |u|+|v|\le L_t:=C[1+q^{-p_*}e^{-b_0w/4}],
 \qquad p_*=603/640.
\]

\begin{proposition}[Value and coefficient budgets for the mass defect]
\label{an03mpv1:massbudgets}
All the following are uniform on $[t_s,t_e]$:
\begin{align}
 \sup|1-M|&\le Cq^{11/10},&
 \int_{t_s}^{t_e}|1-M|&\le C[q^{\alpha_M}+q^{11/10}\ell],
 \label{an03mpv1:massbudget}\\
 \sup |1-M|L_t&\le Cq^{101/640},&
 \int_{t_s}^{t_e}|1-M|L_t
 &\le C[q^{\alpha_M-p_*}+q^{101/640}+q^{11/10}\ell]
   =O(q^{101/640}).
 \label{an03mpv1:linearbudget}
\end{align}
If $Y_*,X_*$ are the selected leading strong comparator coordinates
at $N(t)$, then
\begin{equation}
 \sup |M-1|\bigl(|Y_*|+\Kfun(\rho X_*)\bigr)\le C_*q^{11/10},
 \quad
 \int |M-1|\bigl(|Y_*|+\Kfun(\rho X_*)\bigr)
       \le C_*[q^{\alpha_M}+q^{11/10}\ell].
 \label{an03mpv1:comparator}
\end{equation}
The integrals and suprema in these statements are proved budgets;
no complete TAILREP field-error inequality is asserted here.
\end{proposition}
\begin{proof}
Integrate the two exponentials with their fixed positive rates
$1/3$ and $b_0/4$, and use $t_e-t_s\le20\ell$ for the constant
floor. The smallest power in the product is
$11/10-p_*=101/640>0$, while $\alpha_M\ge199/60$.
This gives \eqref{an03mpv1:massbudget}--\eqref{an03mpv1:linearbudget}.
By Proposition \ref{an03mpv1:phase}, $N(t)<2h<1/2$.
AN06's selected half-learning region bounds $|Y_*|$, $|X_*|$
uniformly, and $\Kfun$ is bounded on that compact interval.
Multiplication gives \eqref{an03mpv1:comparator}.
\end{proof}

\paragraph{Placement in TAILREP's later errors.}
SI/TAILREP's donor comparison retains
$|M-1||Y_*|$ and $|M-1|\Kfun(\rho X_*)$; their additive
supremum budget is $Cq^{11/10}$. P's mass-linear terms are
$-(1-M)v\Phi(u)/2$ and $-(1-M)\Kfun(u)/2$ in the two weak
velocities. Since $0\le\Phi\le1$ and
$\Kfun(u)\le\phi(0)+|u|$, their absolute value on the reached
chart has budget $Cq^{101/640}$. Alternatively, after subtracting
a bounded comparator, their bounded parts cost $Cq^{11/10}$
additively and their linear coefficients can be assigned
$\ell_M(t)=C|1-M(t)|$, with
$\int\ell_M=O(q^{11/10}\ell)$.
These are two possible allocations of the same mass defect, not
two charges to be added. Deriving the remaining original-field terms,
normalized-weight terms and the full common $\ell_B$ inequality
is reserved for S4d.

\section{Step 6: matching theorem, constants and handoff}
\begin{theorem}[Initialized matching package inside the weak passage]
\label{an03mpv1:main}
In the stated deep regime, retain WPT/RB's constants and WBR's
return theorem, fix $a=1/8$, and choose a fixed admissible WBR level
$h$ before choosing $q$. For every sufficiently small $q$ the
original initialized flow on \eqref{an03mpv1:levels} has:
\begin{enumerate}
 \item the full closed clock cohort $S$ in the nonnegative weak-input
 region, with WBR's strict boundary margin and reached chart;
 \item omitted non-core mass $O(q^{1/4}H)$, response
 $c=H(1+O(q^{1/4}))+O(q^2)$, and primary-band amplitude ratio
 $H_{P_q}/H\asymp q^{1/4}$;
 \item the auxiliary logistic comparison \eqref{an03mpv1:logistic}
 with uniform phase \eqref{an03mpv1:phasebound},
 $\int|N-c|=o(1)$ and the interior first-hit relation
 \eqref{an03mpv1:hittime};
 \item exact selected mass identity \eqref{an03mpv1:identity},
 $N_S/H=1+O(q^2)$ and total non-core mass
 $N_{\rm nc}/H=1+O(q^{1/4})$;
 \item for $p=2$ and the actual incoming radial probability,
 $N_0=q^{1/8}$ and $A_0=O(q^{1/16})$;
 \item the pointwise and integrated physical-core mass-defect
 budgets \eqref{an03mpv1:massbudget}--\eqref{an03mpv1:comparator}.
\end{enumerate}
No unproved field-reduction, local strong-reference guard, or
selection hypothesis is imposed to obtain these conclusions.
\end{theorem}
\begin{proof}
WBR supplies return and the chart. Lemma \ref{an03mpv1:band} and
Proposition \ref{an03mpv1:accounting} establish accounting without
using the later sharp transverse conclusion. Proposition
\ref{an03mpv1:phase} then proves logistic matching. BULK's explicitly
matched transverse theorem sharpens the exact mass identity and
proves Proposition \ref{an03mpv1:smallness}; it uses only the
independent signgood energy bound, not logistic matching or smallness.
The last conclusion is Proposition \ref{an03mpv1:massbudgets}.
All intervals are subsets of RB's completed passage.
\end{proof}

\paragraph{Constant order.}
Keep every WPT/RB constant unchanged, then WBR's fixed constants
and boundary margin. The new transverse constant $C_J'$ depends
only on their already fixed $B,L_\times,C_J,b_0$, a fixed upper
bound for $D_q$, and the signgood energy constant. The primary-band
constants depend on fixed Q2EC matching constants and $e^{2D_+}$;
no smallness of $D_+$ is required. Fix the start exponent $a=1/8$
and $p=2$, then take $h$ satisfying WBR's upper bounds, then $q$.
The phase and mass constants have only fixed factors $1/b_0$,
$1/[\lambda(1-h)]$, and the previously fixed passage constants.
No earlier guard is chosen in terms of a new constant. S4d may
decrease $h$ to meet its local small-gain conditions; all statements
here persist for each such fixed choice.

\paragraph{Handoff and scope limit.}
For a later selection endpoint one can take any fixed $n<h$, for
example $n=h/2$, and stop at the auxiliary logistic hit $N=n$.
It lies strictly inside $[t_s,t_e]$, with all actual return and
matching estimates available. The start has the required $p=2$
weighted smallness and no excluded fraction of weak mass. However,
the full weak and strong-reference field inequalities, the local
strong guard $E_0\le\delta/(2K_0)$ (and its actual small-gain
variant), and the selection conclusion remain unproved here.
Fixed original labels do not imply constant normalized radial
weights; a weight-defect coefficient still belongs in S4d's field
audit. The matching theorem never replaces $M$ by one without
\eqref{an03mpv1:massbudget}--\eqref{an03mpv1:comparator}, never
uses the $O(1)$ phase of $t_b$ for an $o(1)$ match, and makes no
continuation claim beyond $t_b$.

\begin{thebibliography}{12}\small
\bibitem{RB} RB: \path{AN03_response_barrier_v1.tex}, prefix
\texttt{an03rbv1:}, \texttt{completed, primaryseed, primaryclock,
clock, conclusions, massconclusions}, S4 handoff.
\bibitem{WBR} WBR: \path{AN03_weak_boundary_return_v1.tex}, prefix
\texttt{an03wbrv1:}, \texttt{main, clockalignment, massrate,
latechart, times}.
\bibitem{SI} SI: \path{AN03_selection_interface_v1.tex}, prefix
\texttt{an03siv1:}, inventory, \texttt{massdeficit, masslinear}.
\bibitem{BULK} BULK: \path{AN03_bulk_clock_and_secondary_Q2_energy_v1.tex},
prefix \texttt{inc:AN03:bulkclock:v1:}, \texttt{bulkclockclosure,
bulkdefectconclusion, forcing, forcinginputs, Jprop, residualinputs}.
\bibitem{RCET} RCET: \path{AN03_residual_closure_and_eta0_transport_v1.tex},
prefix \texttt{an03rcetv1:}, Theorem 2.1 and \texttt{transport}.
\bibitem{LCRC} LCRC: \path{AN03_ledger_cap_residual_closure_v1.tex},
prefix \texttt{an03lcrcv1:}, \texttt{ledger, identity, responsecharge,
signgood}.
\bibitem{Q2EC} Q2EC: \path{AN02_Q2_entry_contracts_v1.tex}, prefix
\texttt{an02q2ecv1:}, \texttt{main, match, coordinatetransport,
transport, targetcontracts}.
\bibitem{Q2EP} Q2EP: \path{AN02_Q2_entry_through_strong_passage_v1.tex},
prefix \texttt{an02q2epv1:}, \texttt{energy, entryenergy}.
\bibitem{TAILREP} TAILREP: \path{AN03_mass_charged_tail_and_selection_defects_v1.tex},
prefix \texttt{inc:AN03:tailrepair:v1:}, \texttt{strongerror,
weakerror, selection}; error placement only, theorem not applied.
\bibitem{WEAK} WEAK: \path{AN03_weak_transit_forcing_and_orbit_selection_v1.tex},
prefix \texttt{inc:AN03:weak:v1:}, \texttt{Wineq}; $W_p$ definition only.
\bibitem{AN06} AN06: \path{AN06_learned_orbit_witnesses_v2.tex},
prefix \texttt{inc:AN06:v2:}, \texttt{regiontheorem}; bounded selected
comparator for $0<N\le1/2$.
\end{thebibliography}
\end{document}
